\documentclass[12 pt]{exam}
\usepackage[margin=0.75in]{geometry}
\usepackage{amsmath,amssymb,color,amsthm, framed}
\usepackage{enumerate,graphicx,multicol}
\usepackage{wrapfig}
\usepackage{pgf,tikz, subcaption}
\usetikzlibrary{arrows}

\newtheorem{claim}{Claim}
\newtheorem{definition}{Definition}
\newcommand{\soln}[1]{\begin{framed}#1\end{framed} \par}

\pagestyle{head}                       % This causes the exam to only have headers, no footers.
%\runningheadrule                     % This causes the headings to have an underline.

%\firstpageheader{Name:\enspace\hbox to 4.5in{\hrulefill} Section:\enspace\hbox to 1.2in{\hrulefill}}{}{}
%\extraheadheight{.25in}

\begin{document}
\hrule
\begin{center}
  \huge{Math 2710 Problem Set 5: \textsection 4.1-5}    %Title
\end{center}
\hrule
\vskip .2in



\begin{questions}
\question Give a direct proof of  the following statement.
If $a$ is an odd integer, then $a+1$ is even.

\question Given component statements $P, Q, R, S,$ and $T$, describe the steps that would need to be used to prove the compound statement $P \wedge Q \implies \sim (R \vee S) \wedge T$ using the method of direct proof.

\question Prove the following. 
Given three sets $A$, $B$, and $C$, $A \subseteq B$ and $A \subseteq C$ if and only if $A \subseteq B \cap C$.

\question Prove that for every set $A$ in a family $\mathcal{S}$ of sets, $\displaystyle{ A \subseteq \bigcup_{S \in \mathcal{S}}S}.$

\question Give a direct proof (or proof by cases) of  the following statement.
For two real numbers $a$ and $b$, $|ab|=|a||b|$.  You may use the following definition:
	\begin{definition}
	The function $|x|$ is defined as $$\begin{cases} x & x \geq 0 \\ -x & x<0 \end{cases}$$
	\end{definition}
	




\section*{Extra Problems}


\question State the result proven below, and determine the proof method used.
	\begin{proof}
		Suppose that $S \subseteq B$ for every $S$ in $\mathcal{S}$, and let $x \in \displaystyle{ \bigcap_{S \in \mathcal{S}}S}$. Then by definition of $\displaystyle{ \bigcap_{S \in \mathcal{S}}S}$, $x$ is an element of some set $T$ in $\mathcal{S}$.
		
		\hspace{.4in} Now, since $S \subseteq B$ for every $S$ in $\mathcal{S}$, it must be true that $T \subseteq B$, and therefore, $x$ is an element of $B$ by definition of $T \subseteq B$. Thus, every element of $\displaystyle{ \bigcap_{S \in \mathcal{S}}S}$ is also an element of $B$.
	\end{proof}
	
\question Review the following proof. Determine the proof method used (or attempted) and give the proof a grade of A (complete, correct, and clear), C (partially correct, partially incomplete, and/or partially unclear), or F (not clear, not correct, or not complete). If you give a grade lower than A, explain your reasoning and describe a way to revise the proof for an A.
	\begin{claim}
		Given a positive real number $x$, the sum of $x$ and its reciprocal is greater than $2$.
	\end{claim}
	\begin{proof}
		We seek to prove that if $x \geq 0$, then $x + \dfrac{1}{x} >2$. If we multiply this inequality by $x$, we get the equivalent inequality $x^2 + 1 \geq 2x$. Subtracting $2x$ from both sides, we see that $x^2 - 2x +1 \geq 0$, or in other words $(x-1)^2 \geq 0$. This is true because any real number's square is always positive or zero. So $x + \dfrac{1}{x} >2$ is true.
	\end{proof}


	

\end{questions}

\end{document}
